
\documentclass{article}
\usepackage{colortbl}
\usepackage{makecell}
\usepackage{multirow}
\usepackage{supertabular}

\begin{document}

\newcounter{utterance}

\centering \large Interaction Transcript for game `cladder', experiment `full\_v1.5\_default', episode 4371 with refineonly{-}Qwen3.5{-}4B.
\vspace{24pt}

{ \footnotesize  \setcounter{utterance}{1}
\setlength{\tabcolsep}{0pt}
\begin{supertabular}{c@{$\;$}|p{.15\linewidth}@{}p{.15\linewidth}p{.15\linewidth}p{.15\linewidth}p{.15\linewidth}p{.15\linewidth}}
    \# & \multicolumn{2}{c}{Player} && \multicolumn{2}{c}{Game Master} \\
    \hline

    \theutterance \stepcounter{utterance}  
    & & & \multicolumn{4}{p{0.6\linewidth}}{
        \cellcolor[rgb]{0.9,0.9,0.9}{
            \makecell[{{p{\linewidth}}}]{
                \texttt{\tiny{[P1$\langle$GM]}}
                \texttt{You are an expert in causal inference. The following question is not a typical commonsense query, but rather a meticulously designed question created by a professor specializing in causal inference, intended to assess the students' mastery of the course content.} \\
\texttt{Imagine a self{-}contained, hypothetical world with only the following conditions, and without any unmentioned factors or causal relationships: Gender has a direct effect on smoking and peanut allergy. Smoking has a direct effect on tar deposit. Tar deposit has a direct effect on peanut allergy. Gender is unobserved. For nonsmokers, the probability of high tar deposit is 28\%. For smokers, the probability of high tar deposit is 58\%. For nonsmokers and with no tar deposit, the probability of being allergic to peanuts is 49\%. For nonsmokers and with high tar deposit, the probability of being allergic to peanuts is 77\%. For smokers and with no tar deposit, the probability of being allergic to peanuts is 27\%. For smokers and with high tar deposit, the probability of being allergic to peanuts is 62\%. The overall probability of smoking is 58\%. Will smoking increase the chance of being allergic to peanuts?} \\
\texttt{Start your answer with 'Yes' or 'No', followed by additional reasoning or evidence to support your explanation.} \\
            }
        }
    }
     \\ \\

    \theutterance \stepcounter{utterance}  
    & \multicolumn{4}{p{0.6\linewidth}}{
        \cellcolor[rgb]{0.561,0.737,0.561}{
            \makecell[{{p{\linewidth}}}]{
                \texttt{\tiny{[P1$\rangle$GM]}}
                \texttt{Yes} \\
\\ 
\texttt{**Reasoning:**} \\
\\ 
\texttt{To determine whether smoking increases the chance of being allergic to peanuts, we must calculate the **causal effect** of smoking on peanut allergy, denoted as \$P(\textbackslash{}text\{Allergy\} | \textbackslash{}text\{do\}(\textbackslash{}text\{Smoker\})) {-} P(\textbackslash{}text\{Allergy\} | \textbackslash{}text\{do\}(\textbackslash{}text\{Nonsmoker\}))\$. This requires isolating the effect of the intervention (forcing a person to smoke or not smoke) from the confounding influence of Gender.} \\
\\ 
\texttt{**1. Define the Causal Structure**} \\
\texttt{Based on the problem description, the causal graph is:} \\
\texttt{*   \$\textbackslash{}text\{Gender\} \textbackslash{}rightarrow \textbackslash{}text\{Smoking\}\$} \\
\texttt{*   \$\textbackslash{}text\{Gender\} \textbackslash{}rightarrow \textbackslash{}text\{Peanut Allergy\}\$} \\
\texttt{*   \$\textbackslash{}text\{Smoking\} \textbackslash{}rightarrow \textbackslash{}text\{Tar Deposit\}\$} \\
\texttt{*   \$\textbackslash{}text\{Tar Deposit\} \textbackslash{}rightarrow \textbackslash{}text\{Peanut Allergy\}\$} \\
\\ 
\texttt{Gender is a confounder because it affects both the exposure (Smoking) and the outcome (Peanut Allergy) directly. Tar Deposit is a mediator.} \\
\\ 
\texttt{**2. Calculate the Counterfactual Probability for Smokers (\$P(\textbackslash{}text\{Allergy\} | \textbackslash{}text\{do\}(\textbackslash{}text\{Smoker\}))\$)**} \\
\texttt{When we intervene to make someone a smoker (\$\textbackslash{}text\{do\}(\textbackslash{}text\{Smoker\})\$), we fix their smoking status regardless of their Gender. We must average the outcomes over the distribution of Gender. However, we do not know the specific probability of Gender (\$P(\textbackslash{}text\{Male\})\$ vs \$P(\textbackslash{}text\{Female\})\$).} \\
\\ 
\texttt{Let \$p = P(\textbackslash{}text\{Gender\} = g)\$ where \$g\$ represents the specific gender group.} \\
\texttt{The probability of being a smoker given gender is \$P(\textbackslash{}text\{Smoker\} | g)\$.} \\
\texttt{The probability of being a nonsmoker given gender is \$P(\textbackslash{}text\{Nonsmoker\} | g) = 1 {-} P(\textbackslash{}text\{Smoker\} | g)\$.} \\
\\ 
\texttt{The overall probability of smoking is given as 0.58:} \\
\texttt{\$\$P(\textbackslash{}text\{Smoker\}) = \textbackslash{}sum\_g P(\textbackslash{}text\{Smoker\} | g) P(g) = 0.58\$\$} \\
\\ 
\texttt{Under the intervention \$\textbackslash{}text\{do\}(\textbackslash{}text\{Smoker\})\$, the probability of having high tar deposit depends only on the fact that they are now smokers:} \\
\texttt{*   \$P(\textbackslash{}text\{High Tar\} | \textbackslash{}text\{do\}(\textbackslash{}text\{Smoker\})) = 0.58\$ (Given in the prompt: "For smokers, the probability of high tar deposit is 58\%").} \\
\texttt{*   \$P(\textbackslash{}text\{High Tar\} | \textbackslash{}text\{do\}(\textbackslash{}text\{Nonsmoker\})) = 0.28\$ (Given in the} \\
            }
        }
    }
    & & \\ \\

    \theutterance \stepcounter{utterance}  
    & & & \multicolumn{2}{p{0.3\linewidth}}{
        \cellcolor[rgb]{0.9,0.9,0.9}{
            \makecell[{{p{\linewidth}}}]{
                \texttt{\tiny{[GM$|$GM]}}
                \texttt{Yes} \\
            }
        }
    }
    & & \\ \\

    \theutterance \stepcounter{utterance}  
    & & & \multicolumn{2}{p{0.3\linewidth}}{
        \cellcolor[rgb]{0.9,0.9,0.9}{
            \makecell[{{p{\linewidth}}}]{
                \texttt{\tiny{[GM$|$GM]}}
                \texttt{yes} \\
            }
        }
    }
    & & \\ \\

    \theutterance \stepcounter{utterance}  
    & & & \multicolumn{2}{p{0.3\linewidth}}{
        \cellcolor[rgb]{0.9,0.9,0.9}{
            \makecell[{{p{\linewidth}}}]{
                \texttt{\tiny{[GM$|$GM]}}
                \texttt{game\_result = WIN} \\
            }
        }
    }
    & & \\ \\

\end{supertabular}
}

\end{document}
